% DOC NO. RL-360-REMOTE-ACCESS · REV. A
%
% This file IS the document. Edit it directly - ripdoc compiles it
% as it stands and stamps the provenance line at build time.
%
%   ripdoc build --only studio-deploy-remote-access
%
% Promoted from studio/deploy/REMOTE-ACCESS.md on 2026-09-07 by `ripdoc promote`.
\documentclass[fontpath=fonts/,logopath=logo/]{riposte-doc}
\usepackage{riposte-md}

\title{Reaching the studio from somewhere else}
\author{Riposte Laboratories}
\date{}
\ripdocno{RL-360-REMOTE-ACCESS}
\riprev{A}
\ripstrapline{The studio is a backend. It accepts uploads and runs \NoCaseChange{\ripmono{ffgac}}/\NoCaseChange{\ripmono{ffedit}} subprocesses on them, and it has \NoCaseChange{\textbf{no authentication of its own}} — on a private network it is gated by the reverse proxy in front of it. That shapes every option below.}
\riprunningtitle{Reaching the studio from somewhere else}

\begin{document}

\ripmakehead

There are two honest ways to use it remotely, and they are very different in what they cost you.


\ripsection{\ripmdhead{1. Over a private overlay (Tailscale) — nothing new to run}}

\textbf{This is already working.} No port is opened, no second copy of the data exists, and nothing is exposed to the internet. Verified 2026-08-14:

\begin{ripmdtable}{X[0.68,l]X[1.32,l]}
\ripmdth{PIECE} & \ripmdth{STATE} \\
Subnet route \ripmono{10.0.1.0/24} & advertised by \ripmono{x}, and it holds the \textbf{primary} route \\
IP forwarding on \ripmono{x} & \ripmono{net.\allowbreak{}ipv4.\allowbreak{}ip\_\allowbreak{}forward =\allowbreak{} 1} \\
Split DNS \ripmono{hq.\allowbreak{}ripostelabs.\allowbreak{}xyz} → \ripmono{10.0.1.99} (Pi-hole) & configured \textbf{tailnet-wide}, with the search domain set too \\
ACL & every device except one gets \ripmono{*:*}, which covers subnet routes \\
TLS & the wildcard \ripmono{*.\allowbreak{}hq.\allowbreak{}ripostelabs.\allowbreak{}xyz} cert already covers the studio's hostname \\
\end{ripmdtable}

So from any tailnet device with \ripmono{-\allowbreak{}-\allowbreak{}accept-\allowbreak{}dns=\allowbreak{}true}, the studio is simply at its normal internal hostname. Nothing about the deployment changes; the Authelia gate still applies, because you are arriving through the same proxy as on the LAN.

\textbf{Adding a new hostname needs no Tailscale change at all.} The split-DNS route covers the whole \ripmono{hq.\allowbreak{}ripostelabs.\allowbreak{}xyz} zone, so a new name resolves the moment Pi-hole and the proxy know about it.


\ripsubsection{\ripmdhead{The one gap: the laptop}}

The tailnet ACL deliberately restricts one device (\ripmono{laptop}) to the mac mini only. It is the single device that \textbf{cannot} reach the studio this way, and it will not fail loudly — it will look like DNS or a dead service.

Tailscale ACLs have no deny rules, so this was implemented by enumerating every \riphi{other} device as an allowed source. Two consequences worth knowing before you debug this at 11pm:

\begin{ripbullets}
\item \textbf{Any newly-added device gets zero access} until it is added to the policy.
\item \textbf{A device removed and re-added can get a new Tailscale IP}, silently breaking the host alias that grants it access.
\end{ripbullets}

Widening this is a deliberate policy decision, not a fix — the restriction was put there on purpose. If you do want the laptop to reach the studio, grant it that destination specifically rather than restoring blanket access.


\ripsubsection{\ripmdhead{Verifying it}}

The pieces above were each checked directly. What was \textbf{not} verified is a full request originating from an off-LAN tailnet device, because no such device on this tailnet accepts a shell (the one that does refuses SSH, and the phone cannot be driven). If you want that proof, from a phone or laptop on the tailnet and off the LAN:

\ripcodeopen
\ripcodesize{9.0}{13.95}
\begin{ripcodeblock}
https://<studio hostname>/api/health     # expect the Authelia login, then ok: true
\end{ripcodeblock}
\ripcodesizereset\ripcodeclose

A 302 to \ripmono{auth.hq…} is the \riphi{correct} result and proves routing and DNS worked. It does not prove the engine is healthy — that is what \ripmono{/api/health} behind the gate is for.


\ripsection{\ripmdhead{2. On a machine of its own}}

Use \ripmono{studio/\allowbreak{}deploy/\allowbreak{}provision.\allowbreak{}sh}. It provisions a fresh Ubuntu 22.04+ or Debian 13 host, x86-64 or arm64, and sizes the unit from that host's own RAM and disk:

\ripcodeopen
\ripcodehead{sh}
\ripcodesize{9.0}{13.95}
\begin{ripcodeblock}
studio/deploy/provision.sh root@<host> --ssh-key ~/.ssh/id_ed25519
\end{ripcodeblock}
\ripcodesizereset\ripcodeclose

It binds to \ripmono{127.0.0.1} by default. Put a gate in front before changing that.


\ripsubsection{\ripmdhead{What it needs}}

\begin{ripmdtable}{X[0.55,l]X[0.55,l]X[0.55,l]X[3.00,l]}
\ripmdth{{}} & \ripmdth{MINIMUM} & \ripmdth{COMFORTABLE} & \ripmdth{WHY} \\
RAM & \textbf{2 GB} & 4 GB & one render peaks near 970 MB (\textasciitilde{}104 B/px × 9M px, measured on \ripmono{colour.hsv}, the costliest op); below a \textasciitilde{}1.4 GB ceiling the cgroup SIGKILLs the engine, which surfaces as a \textbf{502}, not an error \\
CPU & 1 core & 2–4 & \ripmono{GLITZY\_\allowbreak{}MAX\_\allowbreak{}WORKERS} is 1 by default; more workers divide the memory budget, they do not add to it \\
Disk & 10 GB & 20 GB+ & cache budget (default 6 GB) + sources/exports + ffglitch (\textasciitilde{}110 MB) \\
Arch & x86-64 or \textbf{arm64} & {} & ffglitch ships both for Linux \\
\end{ripmdtable}

The script sets \ripmono{GLITZY\_\allowbreak{}MAX\_\allowbreak{}PIXELS}, \ripmono{GLITZY\_\allowbreak{}MAX\_\allowbreak{}WORKERS}, \ripmono{GLITZY\_\allowbreak{}CACHE\_\allowbreak{}BUDGET} and \ripmono{MemoryMax} together from what it finds. They are one setting in four places: \textbf{workers × per-render budget must stay under the cgroup cap.} Raising \ripmono{GLITZY\_\allowbreak{}MAX\_\allowbreak{}PIXELS} alone does not let the engine survive more work — it only stops it refusing work it will then be killed for.


\ripsubsection{\ripmdhead{On Oracle Always Free specifically}}

\begin{ripbullets}
\item \textbf{\ripmono{VM.\allowbreak{}Standard.\allowbreak{}E2.\allowbreak{}1.\allowbreak{}Micro} will not work.} 1 GB of RAM against a 1.4 GB working ceiling; it will be OOM-killed doing ordinary things. The script refuses to provision below 1800 MB rather than let you discover this later.
\item \textbf{\ripmono{VM.\allowbreak{}Standard.\allowbreak{}A1.\allowbreak{}Flex} (Ampere, arm64) is the right shape} — the free allowance is 4 OCPU / 24 GB total and may be \textbf{split across up to 4 instances}, so a 1 OCPU / 6 GB slice is both plenty for the studio and leaves room for other tenants of the same allowance.
\item ffglitch's Linux arm64 build is published as a \textbf{\ripmono{.7z}}, not a \ripmono{.zip}. It is easy to look at the download page, see \ripmono{linux-\allowbreak{}x86\_\allowbreak{}64.\allowbreak{}zip} and \ripmono{macos-\allowbreak{}aarch64.\allowbreak{}zip}, and conclude Linux ARM does not exist. It does, and the provisioner uses it. \ripmono{bsdtar} reads both formats, which is why \ripmono{libarchive-\allowbreak{}tools} is installed rather than p7zip.
\item A1 capacity in single-AD regions is the real obstacle, not the software. There is no second availability domain to fall back to, so a patient retry loop is the only lever — and a \riphi{tight} loop makes it worse, because OCI answers \ripmono{429} and masks the actual capacity error.
\end{ripbullets}


\ripsection{\ripmdhead{Which one}}

If the goal is "use the studio when I'm not at home", \textbf{option 1 is already done} and costs nothing: no new machine, no second copy of the cache, no public surface, and the existing gate keeps applying.

Option 2 is for when the studio needs to serve people who are not on your tailnet — and then the authentication question stops being optional.

\ripend

\end{document}
